\section{A two-arm CRT packet for the minimal affine shear}
\label{sec:affine-two-arm-crt}

Let $a+b=c$ be primitive, put $R=\rad(abc)$, and consider the minimal-step
affine shear
\begin{equation}\label{eq:two-arm-shear}
 U=1+Rh,\qquad V=1+R(h+ck),\qquad W=1+R(h+bk).
\end{equation}
Write $E(n)=n/\rad(n)$.  The density analysis in
Section~\ref{sec:holonomy-density-depth} shows that every three-quarter
exception in the upper-half canonical box must satisfy
\begin{equation}\label{eq:two-arm-marginal-gates}
 Rc<8192E(V),\qquad Rc<8192E(W).
\end{equation}
The two inequalities carry more joint information than two unrelated
numerical lower bounds.

\begin{theorem}[Coprime joint carrier]\label{thm:two-arm-joint-carrier}
If $V,W,T,K$ are positive integers, $(V,W)=1$, and
$T<KE(V)$ and $T<KE(W)$, then
\[
 (E(V),E(W))=1,\qquad E(V)E(W)\mid VW,
 \qquad T^2<K^2E(V)E(W).
\]
\end{theorem}

\begin{proof}
Each powerful part divides its parent integer.  Divisor monotonicity of
coprimality gives $(E(V),E(W))=1$, and multiplication gives the divisibility.
Multiplying the two strict positive inequalities proves the last assertion.
For \eqref{eq:two-arm-marginal-gates}, take $T=Rc$ and $K=8192$.
\end{proof}

The diagonal $h=k$ is automatically admissible because
$(1+Rk,k)=1$, and its long arms are
\[
             V=1+R(c+1)k,\qquad W=1+R(b+1)k.
\]

\begin{theorem}[Diagonal two-arm CRT packet]
\label{thm:two-arm-crt-packet}
Let $D,F$ be positive coprime squarefree integers such that
\[
 (D,R(c+1))=1,\qquad(F,R(b+1))=1.
\]
There is a unique residue class $k_0\pmod {D^2F^2}$ such that
\[
 k\equiv k_0\pmod {D^2F^2}
 \quad\Longleftrightarrow\quad D^2\mid V\ \hbox{ and }\ F^2\mid W.
\]
Every member is admissible and satisfies $D\mid E(V)$ and $F\mid E(W)$.
\end{theorem}

\begin{proof}
The two coefficients $R(c+1)$ and $R(b+1)$ are units modulo $D^2$ and
$F^2$, respectively.  Each affine congruence therefore has one root, and
the Chinese remainder theorem combines them into one class modulo $D^2F^2$.
If a squarefree $D$ satisfies $D^2\mid n$, then each prime of $D$ occurs in
$n$ with valuation at least two and hence occurs in $E(n)$; thus $D\mid E(n)$.
Apply this observation to both arms.  Admissibility was noted above.
\end{proof}

The theorem supplies simultaneous repeated-prime mass, but the following
exact calculation shows that the two marginal gates do not contain enough
information to force exceptionality.

\begin{proposition}[A counted packet and a complete insufficiency witness]
\label{prop:two-arm-counterexample}
For the seed $(a,b,c)=(1,242,243)$ one has $R=66$ and canonical scale
\[
 M=\left\lfloor\frac{243^6}{4\cdot66}\right\rfloor=779890651873.
\]
The diagonal progression
\[
                   h=k=356+1225t
\]
has exactly $318322715$ distinct points with $M/2<k\le M$.  Every point
satisfies $5\mid E(V)$, $7\mid E(W)$, and both inequalities in
\eqref{eq:two-arm-marginal-gates}.  At its first point,
\[
 h=k=389945326231,\qquad (E(U),E(V),E(W))=(1,5,7),
\]
the primitive output is
\[
 (A,B,C)=(25736391531247,
 1519682447137014050,1519708183528545297),\qquad A+B=C,
\]
but $\rad(ABC)^4>C^3$.  Hence the point is not a three-quarter exception.
\end{proposition}

\begin{proof}
On the progression,
\[
 V=25(229321+789096t),\qquad
 W=49(116521+400950t).
\]
The canonical inequalities are equivalent to
$318322715\le t\le636645429$, an interval of exactly $318322715$ integers;
the affine function of $t$ is injective.  The square divisibilities and the
squarefree-carrier lemma give the two excess divisibilities, while
$66\cdot243<8192\cdot5$ and $66\cdot243<8192\cdot7$ give the gates.

For the first point, exact factorization gives
\[
\begin{aligned}
 U&=17\cdot1513905384191,\\
 V&=5^2\cdot23\cdot37\cdot139267\cdot2119433,\\
 W&=7^2\cdot17431\cdot7322098141.
\end{aligned}
\]
All residual factors displayed here are prime.  A shorter certificate for
nonexceptionality is the squarefree divisor
\[
 d=139267\cdot2119433\cdot17431=5145057294975341
\]
of $ABC$: it divides $\rad(ABC)$ and direct integer comparison gives
$d^4>C^3$.
\end{proof}

Proposition~\ref{prop:two-arm-counterexample} refutes only the assertion that
the two necessary long-arm gates are sufficient.  It does not refute their
necessity or the affine route.  The remaining positive theorem must produce,
uniformly over every fixed subcritical seed range $R<c^\lambda$, enough
canonical admissible parameters satisfying the full pointwise inequality
\begin{equation}\label{eq:two-arm-full-gate}
 E(U)E(V)E(W)>
 \frac{RUVW}{(cW)^{3/4}}
 =\frac Rc\,UV(cW)^{1/4}.
\end{equation}
Within the present density architecture, one seeks constants
$\eta,\kappa>0$ and at least
$\kappa R^{-2/3}c^{4+\eta}$ such parameters for every sufficiently large
seed.  Growing squarefree carriers and a correlated powerful-value theorem
for all three arms remain distinct active ways to reach
\eqref{eq:two-arm-full-gate}.
